% ==========================================================
% REFERENCIAS
% Añade cada fuente nueva como un \item más, numerada por
% orden de primera aparición en el texto.
% ==========================================================

\section*{\centering REFERENCIAS}


\begin{enumerate}[leftmargin=1.8em,label=\lbrack\arabic*\rbrack,itemsep=0.6em]

\item F. Ccami-Bernal, P. Urday-Ramos, F. Zela-Coila, J. K. Cáceres-Ruiz y V. Cabrera-Caso, "Trastornos musculoesqueléticos y prácticas ergonómicas en universitarios peruanos durante la pandemia de la COVID-19," \emph{Revista Cubana de Ortopedia y Traumatología}, vol. 37, no. 3, 2023. [En línea]. Disponible en: \url{https://revortopedia.sld.cu/index.php/revortopedia/article/view/608}

\item D. Fernández-Villacorta, A. N. Soriano-Moreno, T. Gálvez-Olórtegui, N. Agüero-Santivañez, D. R. Soriano-Moreno y V. A. Benites-Zapata, "Síndrome visual informático en estudiantes universitarios de posgrado de una universidad privada de Lima, Perú," \emph{Archivos de la Sociedad Española de Oftalmología}, vol. 96, no. 10, 2021. [En línea]. Disponible en: \url{https://www.elsevier.es/es-revista-archivos-sociedad-espanola-oftalmologia-296-articulo-sindrome-visual-informatico-estudiantes-universitarios-S0365669121000058}

\item R. M. Meneses Castañeda, S. L. Ramos Rodríguez, C. D. C. Molfino Jaramillo, E. L. Sánchez Miraval, D. F. Stein Montoros y L. G. Chávez Rodríguez, "Frecuencia del síndrome visual informático en estudiantes de sexto año de medicina durante la educación virtual por COVID-19," \emph{Revista de la Facultad de Medicina Humana}, 2021.

\item Zhunio Suin et al., "Pausas activas y su importancia en el bienestar de estudiantes universitarios," 2024.

\bibitem{5}
R. Hernández-Sampieri, C. Fernández-Collado and P. Baptista-Lucio,
\textit{Metodología de la investigación}, 6th ed.
México D.F., México: McGraw-Hill, 2014.

\item Y. Torres y Y. Rodríguez, ``Surgimiento y evolución de la ergonomía como disciplina: reflexiones sobre la escuela de los factores humanos y la escuela de la ergonomía de la actividad,'' \emph{Revista Facultad Nacional de Salud Pública}, vol. 39, no. 2, e342868, 2021. doi: 10.17533/udea.rfnsp.e342868. [En línea]. Disponible en: \url{https://doi.org/10.17533/udea.rfnsp.e342868}

\item % COMPLETAR autores y año desde la página del artículo
``Factores de riesgo e intervenciones ergonómicas efectivas para el manejo del síndrome de visión de computadora,'' \emph{Ergonomía, Investigación y Desarrollo}, Universidad de Concepción. [En línea]. Disponible en: \url{https://revistas.udec.cl/index.php/Ergonomia_Investigacion/article/download/1345/1974/2396}

\end{enumerate}


